\begin{algorithme}

\procedure{ExplorerForceBrute}
{
    \paramEntree{v : Ville},
    \paramEntree{idCase : \naturelNonNul},
    \paramEntree{idDepart : \naturelNonNul},
    \paramEntree{dirActuelle : Direction},
    \paramEntreeSortie{cheminTemp : \tableauUneDimension{}{de}{\naturel}},
    \paramEntree{profondeur : \naturel},
    \paramEntree{objTrouves : \naturel},
    \paramEntree{objTotal : \naturel},
    \paramEntreeSortie{meilleurChemin : \tableauUneDimension{}{de}{\naturel}},
    \paramEntreeSortie{minLen : \entier},
    \paramEntreeSortie{objCollectes : \tableauUneDimension{}{de}{\booleen}},
    \paramEntree{tailleMax : \naturel}
}
{
    cActuelle : Case \\
    voisinId : \naturel \\
    dir, i : \entier \\
    aRamasseObjet : \booleen
}
{
    \sialors{profondeur < tailleMax \textbf{ et } profondeur < minLen}
    {
        \affecter{cActuelle}{obtenirCase(v, idCase)}

        \sialors{cActuelle \neq \motclefPointeurNULL}
        {
            \affecter{cheminTemp[profondeur]}{idCase}
            \affecter{aRamasseObjet}{faux}


            \sialors{estObligatoire(cActuelle) \textbf{ et } non objCollectes[idCase]}
            {
                \affecter{objCollectes[idCase]}{vrai}
                \affecter{objTrouves}{objTrouves + 1}
                \affecter{aRamasseObjet}{vrai}
            }


            \sialors{idCase = idDepart \textbf{ et } profondeur > 0 \textbf{ et } objTrouves \geq objTotal}
            {
                \affecter{minLen}{profondeur}
                \pour{i}{0}{profondeur}{}{
                    \affecter{meilleurChemin[i]}{cheminTemp[i]}
                }
            }

            \pour{dir}{0}{3}{}{
                \sialors{entierVersDirection(dir) \neq directionOpposee(dirActuelle)}
                {
                    \affecter{voisinId}{obtenirVoisin(cActuelle, entierVersDirection(dir))}
                    \sialors{voisinId \neq 0}
                    {
                        \instruction{ExplorerForceBrute(
                            v,
                            voisinId,
                            idDepart,
                            entierVersDirection(dir),
                            cheminTemp,
                            profondeur + 1,
                            objTrouves,
                            objTotal,
                            meilleurChemin,
                            minLen,
                            objCollectes,
                            tailleMax
                        )}
                    }
                }
            }

            \sialors{aRamasseObjet}
            {
                \affecter{objCollectes[idCase]}{faux}
            }
        }
    }
}


\vspace{1cm}

\begin{algorithme}
\fonction{GenererCheminBrute}
{v : Ville}
{Chemin}
{
    nbCases, departId, objTotal, tailleMax, minLen : Naturel \\
    dirDepart : Direction \\
    cheminTemp, meilleurChemin : \tableauUneDimension{}{de}{Naturel} \\
    objectifs : \tableauUneDimension{}{de}{Naturel} \\
    resultat : Chemin
}
{
    \sialors{v = \motclefPointeurNULL}
    {
        \retourner{creerCheminVide()}
    }

    \affecter{nbCases}{obtenirNbCases(v)}

    \affecter{tailleMax}{(nbCases * 2) + 2}
    \sialors{tailleMax < 4}
    {
        \affecter{tailleMax}{4}
    }

    \affecter{departId}{obtenirNumCaseDepart(v)}
    \affecter{dirDepart}{obtenirDirDepart(v)}
    \affecter{objTotal}{nbCasesObligatoire(v)}

    \allouer{cheminTemp, tailleMax}\\
    \allouer{meilleurChemin, tailleMax}\\
    \allouer{objectifs, nbCases + 1}\\

    \affecter{minLen}{TAILLE\_MAX\_CHEMIN}

    \instruction{ExplorerForceBrute(
        v,
        departId,
        departId,
        dirDepart,
        cheminTemp,
        0,
        0,
        objTotal,
        meilleurChemin,
        minLen,
        objectifs,
        tailleMax
    )}

    \affecter{resultat}{construireCheminFinal(
        v,
        meilleurChemin,
        minLen,
        departId,
        dirDepart
    )}

    \desallouer{cheminTemp}\\
    \desallouer{meilleurChemin}\\
    \desallouer{objectifs}\\

    \retourner{resultat}
}
\end{algorithme}